% ==============================================================================
% Elsevier journal submission template (skeleton)
% Based on elsarticle v3.5  —  CTAN: https://ctan.org/pkg/elsarticle
% Class elsarticle is contained in standard TeX Live / MiKTeX, so this file
% compiles out-of-the-box with:  pdflatex -> bibtex -> pdflatex -> pdflatex
%
% Class modes (pick exactly one):
%   1p   single-column, preprint style (loosest)
%   3p   single-column, review style with line numbers  (DEFAULT for review)
%   5p   two-column, final typeset style (closest to published article)
% Class options:
%   times           use Times font (recommended)
%   review          double-spacing + line numbers (review copy)
%   authoryear      author-year citation
%   number          numbered citation
%   nonatbib        do not load natbib
%
% Typography (elsarticle defaults, with `times` option):
%   Body font : Times Roman 10pt
%   Math font : Times (mathptmx)
%   Spacing   : single (review: double)
% Citation   : natbib, [1] or (Author, Year) depending on options
% Bibliography:
%   elsarticle-num         numbered, citation order
%   elsarticle-num-names   numbered, with author names
%   elsarticle-harv        author-year (Harvard)
% Figures    : single column = \columnwidth (3.5in) | double = \textwidth
% Captions   : below figures, above tables; 9pt
% ==============================================================================
\documentclass[3p,times]{elsarticle}
% For a two-column final look:
% \documentclass[5p,times]{elsarticle}
% For a preprint look:
% \documentclass[1p,times]{elsarticle}
% For author-year citations, use \documentclass[3p,times]{elsarticle} with
%   \usepackage[authoryear]{natbib}  (requires the `nonatbib` class option)

% ---- Core packages (natbib is auto-loaded by elsarticle; lineno is too) ----
\usepackage{amsmath,amssymb,amsfonts}
\usepackage{graphicx}
\usepackage{booktabs}
\usepackage{array}
\usepackage{xcolor}
\usepackage{algorithm}
\usepackage{algorithmic}
% \usepackage{hyperref}          % load last if used
% \hypersetup{colorlinks=true,linkcolor=blue,citecolor=blue,urlcolor=blue}

% \modulolinenumbers[1]          % uncomment for line numbers (review mode)

% ==============================================================================
% Frontmatter: title, authors, affiliations, abstract, keywords
% ==============================================================================
\begin{document}

\begin{frontmatter}

\title{A Concise and Informative Title for the Elsevier Submission}

\author[a]{First Author\corref{cor1}}
\ead{first.author@univ.edu}
\cortext[cor1]{Corresponding author.}

\author[b]{Second Author}

\author[a]{Senior Author}

\address[a]{Department of X, University of Y, City, Country}
\address[b]{Department of Z, Institute of W, City, Country}

\begin{abstract}
% TODO: a single structured abstract (typically 150-250 words). State the
% problem, approach, key results with effect sizes, and the conclusion.
% Elsevier journals often prefer a structured abstract; check the journal's
% guide for authors. No citations in the abstract.
\end{abstract}

\begin{keyword}
% TODO: 4-6 keywords separated by `;`  (Elsevier convention).
keyword one \sep keyword two \sep keyword three \sep keyword four
\end{keyword}

\end{frontmatter}

% ==============================================================================
% 1. Introduction
% ==============================================================================
\section{Introduction}
\label{sec:intro}
% TODO: context -> gap -> approach -> contributions -> roadmap. Use \citep{}
% for parenthetical citations and \citet{} for textual citations (natbib).

% ==============================================================================
% 2. Related Work
% ==============================================================================
\section{Related Work}
\label{sec:related}
% TODO: cluster prior work; contrast methodology, data, and findings;
% explicitly state how this work differs.

% ==============================================================================
% 3. Methods
% ==============================================================================
\section{Methods}
\label{sec:methods}
% TODO: problem formulation, notation, proposed method, theoretical analysis.

\subsection{Problem formulation}
% TODO: state the problem formally; define notation in a table if heavy.

\subsection{Proposed method}
% TODO: describe the method; include pseudo-code via the algorithm environment.

\begin{algorithm}[t]
\caption{Outline of the proposed method}
\label{alg:method}
\begin{algorithmic}[1]
\STATE \textbf{Input:} dataset $\mathcal{D}$, hyper-parameters $\Theta$
\STATE initialise $\theta$
\WHILE{not converged}
    \STATE $\theta \leftarrow \theta - \eta\,\nabla\mathcal{L}(\theta;\mathcal{D})$
\ENDWHILE
\STATE \textbf{return} $\theta$
\end{algorithmic}
\end{algorithm}

% ==============================================================================
% 4. Experiments
% ==============================================================================
\section{Experiments}
\label{sec:experiments}
% TODO: setup (datasets / baselines / metrics / hardware), main results,
% ablation, qualitative examples, statistical significance.

\begin{figure}[t]
\centering
% \includegraphics[width=\columnwidth]{figures/result.pdf}
\fbox{\rule{0pt}{4cm}\rule{\columnwidth}{0pt}}
\caption{Single-column figure. Caption goes BELOW the figure. Use
\texttt{\textbackslash columnwidth} for single-column, \texttt{\textbackslash textwidth}
or \texttt{figure*} for double-column.}
\label{fig:result}
\end{figure}

\begin{table}[t]
\caption{Table caption goes ABOVE the table (Elsevier convention).}
\label{tab:results}
\centering
\begin{tabular}{lcc}
\toprule
Method   & Metric A       & Metric B       \\
\midrule
Baseline & 0.812          & 0.754          \\
Ours     & \textbf{0.891} & \textbf{0.832} \\
\bottomrule
\end{tabular}
\end{table}

% ==============================================================================
% 5. Conclusion
% ==============================================================================
\section{Conclusion}
\label{sec:conclusion}
% TODO: recap contribution and headline result; one sentence of future work.

% ==============================================================================
% Acknowledgements
% ==============================================================================
\section*{Acknowledgements}
% TODO: funding, grant numbers, individuals.

% ==============================================================================
% References  —  pick the .bst that matches your class options:
%   elsarticle-num         (numbered)            for [1p/3p/5p, number]
%   elsarticle-num-names   (numbered + names)
%   elsarticle-harv        (author-year)         for [*, authoryear]
% Compile:  pdflatex main  ->  bibtex main  ->  pdflatex main  ->  pdflatex main
% ==============================================================================
\bibliographystyle{elsarticle-num}
\bibliography{references}
% Inline fallback (uncomment if not using BibTeX):
% \begin{thebibliography}{99}
% \bibitem{ref_example} F.~Author and S.~Author, Title of the paper,
%   \emph{Journal Name} XX (20YY) 1--10.
% \end{thebibliography}

\end{document}
